
\section{Limitations and Future Work}
\label{Limitation and future work}

While SECDA-DSE demonstrates promising early-stage results across multiple accelerator workloads, the framework is still in an active development phase and currently presents some limitations. 
One of the primary challenges is that a sufficiently diverse set of hardware datapoints are initially required for effective fine-tuning and refinement of the LLM Stack. 
Furthermore, the current workflow still incorporates a human-in-the-loop evaluation stage for validating generated designs before FPGA execution. Although this improves reliability during early-stage development, scaling the framework toward larger workload coverage and broader exploration spaces will require significantly higher levels of automation. 
Another limitation is that the current evaluation covers only a small set of workloads on a single FPGA platform. Although the generated accelerators demonstrate workload-aware behavior and successful FPGA execution, broader evaluations across more workloads, architectures, and FPGA families are needed to assess the generalization capability of SECDA-DSE.

In future work, our goal is to fully automate the accelerator evaluation workflow by integrating the SECDA toolkit directly into SECDA-DSE through a model context protocol (MCP) server. This integration will allow the LLM Stack to directly invoke SECDA testing, simulation, HLS, and hardware execution functions without requiring manual intervention, enabling a more autonomous hardware exploration workflow. 
Additionally, the current evaluation focuses on a relatively small set of workloads and a single FPGA platform, as future work we plan to extend evaluation of SECDA-DSE across multiple parameter scales, LLMs and FPGAs to evaluate how model size, architecture and choice of the FPGA board influence accelerator generation quality and hardware exploration capability.